% Options for packages loaded elsewhere
\PassOptionsToPackage{unicode}{hyperref}
\PassOptionsToPackage{hyphens}{url}
\documentclass[
  11pt,
]{article}
\usepackage{xcolor}
\usepackage[margin=1in]{geometry}
\usepackage{amsmath,amssymb}
\setcounter{secnumdepth}{-\maxdimen} % remove section numbering
\usepackage{iftex}
\ifPDFTeX
  \usepackage[T1]{fontenc}
  \usepackage[utf8]{inputenc}
  \usepackage{textcomp} % provide euro and other symbols
\else % if luatex or xetex
  \usepackage{unicode-math} % this also loads fontspec
  \defaultfontfeatures{Scale=MatchLowercase}
  \defaultfontfeatures[\rmfamily]{Ligatures=TeX,Scale=1}
\fi
\usepackage{lmodern}
\ifPDFTeX\else
  % xetex/luatex font selection
\fi
% Use upquote if available, for straight quotes in verbatim environments
\IfFileExists{upquote.sty}{\usepackage{upquote}}{}
\IfFileExists{microtype.sty}{% use microtype if available
  \usepackage[]{microtype}
  \UseMicrotypeSet[protrusion]{basicmath} % disable protrusion for tt fonts
}{}
\makeatletter
\@ifundefined{KOMAClassName}{% if non-KOMA class
  \IfFileExists{parskip.sty}{%
    \usepackage{parskip}
  }{% else
    \setlength{\parindent}{0pt}
    \setlength{\parskip}{6pt plus 2pt minus 1pt}}
}{% if KOMA class
  \KOMAoptions{parskip=half}}
\makeatother
\usepackage{color}
\usepackage{fancyvrb}
\newcommand{\VerbBar}{|}
\newcommand{\VERB}{\Verb[commandchars=\\\{\}]}
\DefineVerbatimEnvironment{Highlighting}{Verbatim}{commandchars=\\\{\}}
% Add ',fontsize=\small' for more characters per line
\newenvironment{Shaded}{}{}
\newcommand{\AlertTok}[1]{\textcolor[rgb]{1.00,0.00,0.00}{\textbf{#1}}}
\newcommand{\AnnotationTok}[1]{\textcolor[rgb]{0.38,0.63,0.69}{\textbf{\textit{#1}}}}
\newcommand{\AttributeTok}[1]{\textcolor[rgb]{0.49,0.56,0.16}{#1}}
\newcommand{\BaseNTok}[1]{\textcolor[rgb]{0.25,0.63,0.44}{#1}}
\newcommand{\BuiltInTok}[1]{\textcolor[rgb]{0.00,0.50,0.00}{#1}}
\newcommand{\CharTok}[1]{\textcolor[rgb]{0.25,0.44,0.63}{#1}}
\newcommand{\CommentTok}[1]{\textcolor[rgb]{0.38,0.63,0.69}{\textit{#1}}}
\newcommand{\CommentVarTok}[1]{\textcolor[rgb]{0.38,0.63,0.69}{\textbf{\textit{#1}}}}
\newcommand{\ConstantTok}[1]{\textcolor[rgb]{0.53,0.00,0.00}{#1}}
\newcommand{\ControlFlowTok}[1]{\textcolor[rgb]{0.00,0.44,0.13}{\textbf{#1}}}
\newcommand{\DataTypeTok}[1]{\textcolor[rgb]{0.56,0.13,0.00}{#1}}
\newcommand{\DecValTok}[1]{\textcolor[rgb]{0.25,0.63,0.44}{#1}}
\newcommand{\DocumentationTok}[1]{\textcolor[rgb]{0.73,0.13,0.13}{\textit{#1}}}
\newcommand{\ErrorTok}[1]{\textcolor[rgb]{1.00,0.00,0.00}{\textbf{#1}}}
\newcommand{\ExtensionTok}[1]{#1}
\newcommand{\FloatTok}[1]{\textcolor[rgb]{0.25,0.63,0.44}{#1}}
\newcommand{\FunctionTok}[1]{\textcolor[rgb]{0.02,0.16,0.49}{#1}}
\newcommand{\ImportTok}[1]{\textcolor[rgb]{0.00,0.50,0.00}{\textbf{#1}}}
\newcommand{\InformationTok}[1]{\textcolor[rgb]{0.38,0.63,0.69}{\textbf{\textit{#1}}}}
\newcommand{\KeywordTok}[1]{\textcolor[rgb]{0.00,0.44,0.13}{\textbf{#1}}}
\newcommand{\NormalTok}[1]{#1}
\newcommand{\OperatorTok}[1]{\textcolor[rgb]{0.40,0.40,0.40}{#1}}
\newcommand{\OtherTok}[1]{\textcolor[rgb]{0.00,0.44,0.13}{#1}}
\newcommand{\PreprocessorTok}[1]{\textcolor[rgb]{0.74,0.48,0.00}{#1}}
\newcommand{\RegionMarkerTok}[1]{#1}
\newcommand{\SpecialCharTok}[1]{\textcolor[rgb]{0.25,0.44,0.63}{#1}}
\newcommand{\SpecialStringTok}[1]{\textcolor[rgb]{0.73,0.40,0.53}{#1}}
\newcommand{\StringTok}[1]{\textcolor[rgb]{0.25,0.44,0.63}{#1}}
\newcommand{\VariableTok}[1]{\textcolor[rgb]{0.10,0.09,0.49}{#1}}
\newcommand{\VerbatimStringTok}[1]{\textcolor[rgb]{0.25,0.44,0.63}{#1}}
\newcommand{\WarningTok}[1]{\textcolor[rgb]{0.38,0.63,0.69}{\textbf{\textit{#1}}}}
\usepackage{longtable,booktabs,array}
\newcounter{none} % for unnumbered tables
\usepackage{calc} % for calculating minipage widths
% Correct order of tables after \paragraph or \subparagraph
\usepackage{etoolbox}
\makeatletter
\patchcmd\longtable{\par}{\if@noskipsec\mbox{}\fi\par}{}{}
\makeatother
% Allow footnotes in longtable head/foot
\IfFileExists{footnotehyper.sty}{\usepackage{footnotehyper}}{\usepackage{footnote}}
\makesavenoteenv{longtable}
\usepackage{graphicx}
\makeatletter
\newsavebox\pandoc@box
\newcommand*\pandocbounded[1]{% scales image to fit in text height/width
  \sbox\pandoc@box{#1}%
  \Gscale@div\@tempa{\textheight}{\dimexpr\ht\pandoc@box+\dp\pandoc@box\relax}%
  \Gscale@div\@tempb{\linewidth}{\wd\pandoc@box}%
  \ifdim\@tempb\p@<\@tempa\p@\let\@tempa\@tempb\fi% select the smaller of both
  \ifdim\@tempa\p@<\p@\scalebox{\@tempa}{\usebox\pandoc@box}%
  \else\usebox{\pandoc@box}%
  \fi%
}
% Set default figure placement to htbp
\def\fps@figure{htbp}
\makeatother
\setlength{\emergencystretch}{3em} % prevent overfull lines
\providecommand{\tightlist}{%
  \setlength{\itemsep}{0pt}\setlength{\parskip}{0pt}}
\usepackage{bookmark}
\IfFileExists{xurl.sty}{\usepackage{xurl}}{} % add URL line breaks if available
\urlstyle{same}
\hypersetup{
  hidelinks,
  pdfcreator={LaTeX via pandoc}}

\author{}
\date{}


\usepackage{seqsplit}
\tolerance=9999
\emergencystretch=3em
\hbadness=9999
\begin{document}

\section{Topological Mismatch in Distributed Inference: Mechanized
Models and Protocol-Level Evidence for Fork/Race/Fold
Scheduling}\label{topological-mismatch-in-distributed-inference-mechanized-models-and-protocol-level-evidence-for-forkracefold-scheduling}

\emph{\href{https://www.patreon.com/cw/twbuley}{Taylor William Buley} --
Independent Researcher}\\
\emph{\href{mailto:taylor@forkjoin.ai}{\nolinkurl{taylor@forkjoin.ai}}}

\subsection{Abstract}\label{abstract}

Path-like execution can impose avoidable coordination cost when the
workload's natural structure has positive parallel-path complexity. This
paper studies that mismatch in a finite-DAG fork/race/fold model using
the diagnostic Delta\_beta = beta\_1* - beta\_1\^{}exec, operationalized
by bounded graph extraction at a fixed layer, and a correctness envelope
C1-C4. Within that modeled scope, the paper proves a pipeline speedup
floor and sandwich, classifies Reynolds regimes by Re = N / C, and
separates transport path-collapse boundaries from within-class protocol
efficiency. The engineering surface instantiates the model in
distributed staged computation and Aeon Flow transport. In the companion
artifacts, 8/8 primary wall-clock cells pass with median speedups
11.785x-21.620x, and an Aeon Flow versus HTTP/3 protocol corpus over 144
sites and 12,371 resources shows a 72.252\% median framing gain with
positive latency gains across all 6 primary environments. The claim is
not that fork/race/fold is universally optimal; it is that bounded graph
pairs make Delta\_beta falsifiable and useful for diagnosing when
path-like execution is misfit to the workload.

\subsection{1. Problem and Claim}\label{problem-and-claim}

\textbf{Thesis.} Path-like execution topologies impose avoidable
coordination cost when the workload's natural structure has positive
parallel-path complexity.

The narrow target in this paper is distributed inference and the
protocol stack around it. A staged inference workload naturally contains
branch points, joins, and independent paths that can be processed
without violating correctness. A path-like runtime collapses that
structure into a single ordered line. Sometimes that is exactly the
right choice. Sometimes it is a mismatch that inserts waiting, framing
waste, or unnecessary synchronization even though the workload itself
would allow more structure.

This paper isolates that mismatch as a topological question rather than
a hardware or implementation-style question. The central diagnostic is
\texttt{Delta\_beta}, the gap between the parallel-path complexity
required by the workload and the path budget actually exposed by the
execution topology. When \texttt{Delta\_beta\ =\ 0}, the schedule is
structurally matched to the workload at the chosen abstraction layer.
When \texttt{Delta\_beta\ \textgreater{}\ 0}, some coordination cost is
being imposed by the topology rather than by the problem.

\textbf{What is new here.} The flagship claim is not that overlap,
pipelining, or multi-stream transport were previously unknown. The
novelty is narrower.

\begin{enumerate}
\def\labelenumi{\arabic{enumi}.}
\tightlist
\item
  It separates workload-side path demand \texttt{beta\_1*} from
  runtime-side path budget \texttt{beta\_1\^{}exec}, computes both from
  bounded graph pairs at the same declared layer, and uses
  \texttt{Delta\_beta} as the mismatch coordinate between them.
\item
  It binds the main body claims to a theorem-to-observable surface: each
  positive claim has a local formal anchor, a measured observable, and a
  stated falsifier.
\item
  It treats scheduler and transport as one coupled execution surface, so
  the wire is allowed to reintroduce path-like collapse even when the
  scheduler itself exposes concurrency.
\end{enumerate}

The contribution is intentionally narrow.

\begin{enumerate}
\def\labelenumi{\arabic{enumi}.}
\tightlist
\item
  It gives a finite DAG scheduling model with four primitives: fork,
  race, fold, and vent.
\item
  It states an explicit correctness envelope, C1-C4, under which the
  model is used.
\item
  It identifies a small theorem-to-observable bridge sufficient for the
  scheduler and protocol claims used here.
\item
  It evaluates that bridge with artifact-backed wall-clock evidence,
  protocol-corpus evidence, and an explicit path-boundary control rather
  than with analogy-driven extrapolation.
\end{enumerate}

The flagship stays short on purpose because the claim is local and
falsifiable. If the bridge needs a survey-length manuscript to survive
compression, it is not yet sharp enough.

The point is not to rename ordinary pipelining. Classical pipelining
already explains why overlapping stages can reduce latency. Dataflow and
bounded DAG systems already explain how forks and joins can be made
explicit. Modern transports already reduce some head-of-line and framing
costs. The stronger claim here is different: some execution failures are
best understood as forcing a workload with positive intrinsic
parallel-path complexity through a topology that exposes too little of
that structure. In that setting, the useful question is not only ``how
much parallel hardware exists?'' but also ``how much of the workload's
path structure has the runtime allowed to exist, and where is the
missing path budget being collapsed?''

The scheduler and transport story are studied together because they fail
together. A staged pipeline can preserve local concurrency and still
lose it at the protocol layer if the wire format, ordering policy, or
collapse semantics reserialize the flow. Likewise, a richer transport
can still fail to help if the scheduler never exposes multiple useful
paths. The flagship claim is therefore a bridge claim: the same mismatch
diagnostic can be used to reason about staged execution and about the
protocol that carries staged execution across machines.

This paper does not ask the reader to grant a universal theory. It asks
for a smaller standard. Given a finite DAG workload under explicit
correctness conditions, do the mechanized model and the checked
artifacts jointly support the claim that path-like execution can impose
avoidable coordination cost? The answer argued here is yes.

\subsection{2. Model and Notation}\label{model-and-notation}

\subsubsection{2.1 Finite DAG Scheduling
Surface}\label{finite-dag-scheduling-surface}

Let \texttt{G\ =\ (V,\ E)} be a finite DAG representing a bounded
execution slice. Nodes denote local computations or transport actions.
Edges denote causality. A fork node has out-degree greater than one, a
fold node has in-degree greater than one, and a linear chain has
in-degree and out-degree at most one. The model uses four operational
primitives:

\begin{enumerate}
\def\labelenumi{\arabic{enumi}.}
\tightlist
\item
  \texttt{Fork(S,\ O)\ -\textgreater{}\ \{S\_1,\ ...,\ S\_k\}} creates
  branch-local states.
\item
  \texttt{Race(\{S\_i\})\ -\textgreater{}\ (S\_w,\ i\_w)} advances
  branches concurrently and accepts the first valid completion.
\item
  \texttt{Fold(\{S\_i\},\ f)\ -\textgreater{}\ S*} merges branch
  outcomes with a deterministic fold function \texttt{f}.
\item
  \texttt{Vent(S\_i)\ -\textgreater{}\ bottom} terminates a branch
  without corrupting its siblings.
\end{enumerate}

The paper does not need a larger primitive set. The local questions are
scheduling questions: when is a path split useful, when does a branch
need an all-branch fold, and when can an execution surface safely accept
early completion or local failure?

\subsubsection{2.2 Parallel-Path Complexity and
Delta\_beta}\label{parallel-path-complexity-and-delta_beta}

Let \texttt{beta\_1*} denote the intrinsic parallel-path complexity of
the workload at the abstraction layer under study. This is not claimed
to be an invariant of reality at every scale; it is a model-side count
of how many independent paths the workload can profitably expose while
preserving correctness. Let \texttt{beta\_1\^{}exec} denote the path
budget exposed by the execution topology at that same layer. The
diagnostic is

\[
\Delta_\beta = \beta_1^* - \beta_1^{\mathrm{exec}}.
\]

If \texttt{Delta\_beta\ =\ 0}, the schedule is structurally matched. If
\texttt{Delta\_beta\ \textgreater{}\ 0}, the execution surface is
forcing some part of the workload through a narrower path than the
workload requires. The diagnostic is directional and local. It is not a
universal score. It only answers whether the chosen runtime is narrower
than the modeled workload.

To keep that claim falsifiable, the paper uses one explicit extraction
rule. Fix one abstraction layer and one bounded execution slice. Build a
workload graph \texttt{G*} whose nodes are bounded events and whose
edges are only the dependencies required for correctness at that layer.
Build an effective execution graph \texttt{G\^{}exec} for the same slice
after scheduler and transport ordering have been applied. Then compute

\[
\beta_1(G) = \max(0, |E| - |V| + \beta_0(G))
\]

on each graph, using the same event vocabulary on both sides. In the
companion, this is implemented directly over bounded trace and protocol
graphs rather than inferred from benchmark scalars alone.

A concrete worked example is enough to show the rule. A 10-stream
fork/join workload has \texttt{\textbar{}V\textbar{}\ =\ 12},
\texttt{\textbar{}E\textbar{}\ =\ 20}, and \texttt{beta\_0\ =\ 1}, so
\texttt{beta\_1*\ =\ 9}. If an ordered single-pipe transport forces the
same slice into a chain, then \texttt{\textbar{}E\textbar{}\ =\ 11},
\texttt{beta\_1\^{}exec\ =\ 0}, and \texttt{Delta\_beta\ =\ 9}. If the
effective transport preserves the same 10 branch-local streams through
independent recovery, then \texttt{beta\_1\^{}exec\ =\ 9} and
\texttt{Delta\_beta\ =\ 0}. This paper uses that graph-pair rule as the
operational meaning of the diagnostic.

Three execution classes matter in this paper.

{\def\LTcaptype{none} % do not increment counter
\begin{longtable}[]{@{}
  >{\raggedright\arraybackslash}p{(\linewidth - 6\tabcolsep) * \real{0.2308}}
  >{\raggedright\arraybackslash}p{(\linewidth - 6\tabcolsep) * \real{0.2308}}
  >{\raggedleft\arraybackslash}p{(\linewidth - 6\tabcolsep) * \real{0.3077}}
  >{\raggedright\arraybackslash}p{(\linewidth - 6\tabcolsep) * \real{0.2308}}@{}}
\toprule\noalign{}
\begin{minipage}[b]{\linewidth}\raggedright
Execution class
\end{minipage} & \begin{minipage}[b]{\linewidth}\raggedright
Informal topology
\end{minipage} & \begin{minipage}[b]{\linewidth}\raggedleft
Typical \texttt{beta\_1\^{}exec} reading
\end{minipage} & \begin{minipage}[b]{\linewidth}\raggedright
Use in this paper
\end{minipage} \\
\midrule\noalign{}
\endhead
\bottomrule\noalign{}
\endlastfoot
Sequential path & one ordered line & \texttt{0} & boundary case \\
Chunked pipeline & stage-local concurrency without workload branching &
positive but bounded by stage occupancy & staged baseline \\
Fork/race/fold schedule & explicit branch, race, and deterministic
collapse surface & matched to positive \texttt{beta\_1*} when the
workload contains multiple useful paths & target topology \\
\end{longtable}
}

The flagship claim is not that every problem has positive
\texttt{beta\_1*}. Some workloads are truly path-like. For those,
\texttt{Delta\_beta\ =\ 0} under a sequential schedule and no theorem in
this paper predicts a win from adding concurrency. The diagnostic
becomes useful only when the workload is naturally richer than the
schedule that carries it.

The paper does not ask \texttt{Delta\_beta} to replace classical
metrics. It adds it because the existing metrics answer different
questions.

{\def\LTcaptype{none} % do not increment counter
\begin{longtable}[]{@{}
  >{\raggedright\arraybackslash}p{(\linewidth - 6\tabcolsep) * \real{0.2500}}
  >{\raggedright\arraybackslash}p{(\linewidth - 6\tabcolsep) * \real{0.2500}}
  >{\raggedright\arraybackslash}p{(\linewidth - 6\tabcolsep) * \real{0.2500}}
  >{\raggedright\arraybackslash}p{(\linewidth - 6\tabcolsep) * \real{0.2500}}@{}}
\toprule\noalign{}
\begin{minipage}[b]{\linewidth}\raggedright
Quantity
\end{minipage} & \begin{minipage}[b]{\linewidth}\raggedright
What it answers
\end{minipage} & \begin{minipage}[b]{\linewidth}\raggedright
What it is useful for
\end{minipage} & \begin{minipage}[b]{\linewidth}\raggedright
What it does not localize
\end{minipage} \\
\midrule\noalign{}
\endhead
\bottomrule\noalign{}
\endlastfoot
serial fraction & how much work is currently marked serial & classical
scaling ceilings & whether the serialization is inherent to the workload
or inserted by the runtime \\
occupancy / utilization & how full stages or streams are & idle-slot and
saturation diagnosis & whether a lower layer has silently collapsed the
path structure back to one ordered flow \\
exposed path budget \texttt{beta\_1\^{}exec} & how many independent
paths the runtime actually keeps alive & describing the execution
surface & whether the workload needed more than that many paths \\
\texttt{Delta\_beta} & how much modeled path structure is missing from
the execution surface & locating topological mismatch & absolute system
quality outside the chosen model boundary \\
\end{longtable}
}

The point is therefore comparative, not imperial. Serial fraction and
occupancy remain useful. \texttt{Delta\_beta} earns its keep when they
are not enough. An application can expose multiple logical streams and
still inherit a path-like transport below it. In that case occupancy at
the stream layer may look healthy while the transport still forces
\texttt{beta\_1\^{}exec\ =\ 0}. The diagnostic names that contradiction
directly. Conversely, if the workload is genuinely one-path, then
\texttt{beta\_1*\ =\ 0}, \texttt{Delta\_beta\ =\ 0}, and the model
predicts no structural win from richer scheduling.

That is the indispensable use-case for \texttt{Delta\_beta} in this
paper: concurrency can be visible in the application graph and absent in
the effective execution graph at the same time.

\subsubsection{2.3 C1-C4 Correctness
Envelope}\label{c1-c4-correctness-envelope}

The finite DAG model is used only inside the following correctness
envelope.

{\def\LTcaptype{none} % do not increment counter
\begin{longtable}[]{@{}
  >{\raggedright\arraybackslash}p{(\linewidth - 4\tabcolsep) * \real{0.3333}}
  >{\raggedright\arraybackslash}p{(\linewidth - 4\tabcolsep) * \real{0.3333}}
  >{\raggedright\arraybackslash}p{(\linewidth - 4\tabcolsep) * \real{0.3333}}@{}}
\toprule\noalign{}
\begin{minipage}[b]{\linewidth}\raggedright
Condition
\end{minipage} & \begin{minipage}[b]{\linewidth}\raggedright
Statement
\end{minipage} & \begin{minipage}[b]{\linewidth}\raggedright
Operational reading
\end{minipage} \\
\midrule\noalign{}
\endhead
\bottomrule\noalign{}
\endlastfoot
C1 & Constraint locality & Stage-local ordering is sufficient for global
correctness. \\
C2 & Branch isolation & A vented branch cannot corrupt sibling
branches. \\
C3 & Deterministic fold & The merge function is deterministic on the
admitted branch outputs. \\
C4 & Termination & Every branch completes, vents, or times out in finite
time. \\
\end{longtable}
}

These are not cosmetic assumptions. If C1 fails, the system may require
a stronger global order than the schedule preserves. If C2 fails, a
local rejection can leak damage sideways. If C3 fails, fold introduces
ambiguity rather than removing it. If C4 fails, the path budget never
collapses and the schedule cannot be treated as a bounded computation.

Within the formal companion, C1-C4 are checked in two ways. TLA+ models
make them explicit as state invariants or progress constraints on
bounded executions. Lean 4 theorem schemas supply the quantitative
consequences that depend on those conditions. The paper uses only the
part of that surface needed for the pipeline and protocol claims below.

\subsubsection{2.4 Pipeline Variables}\label{pipeline-variables}

For a prompt or batch with \texttt{P} items, chunk size \texttt{B},
stage count \texttt{N}, and chunk count

\[
C = \lceil P / B \rceil,
\]

the chunked pipeline time under the modeled schedule is

\[
T_{\mathrm{pipe}} = \lceil P/B \rceil + N - 1,
\]

while the fully serialized baseline is

\[
T_{\mathrm{seq}} = P \cdot N.
\]

The step-count speedup under the stated idealization is therefore

\[
\mathrm{speedup} = \frac{P \cdot N}{\lceil P/B \rceil + N - 1}.
\]

The model uses two explicit assumptions for this equation.

\begin{enumerate}
\def\labelenumi{\arabic{enumi}.}
\tightlist
\item
  A1: per-chunk stage service times are homogeneous across stages.
\item
  A2: inter-stage communication and synchronization cost is zero.
\end{enumerate}

Those assumptions matter. The step-count theorem is a model claim about
scheduling shape, not a promise that real deployments will achieve the
same scalar under arbitrary transport, compute heterogeneity, or loss.

\subsubsection{2.5 Reynolds Regimes}\label{reynolds-regimes}

The pipeline Reynolds number is defined as

\[
Re = N / C.
\]

This quantity measures how dense the pipeline is in stages per chunk.
Low \texttt{Re} means many chunks relative to stage count, which pushes
the system toward steady-state occupancy. High \texttt{Re} means few
chunks relative to stage count, which magnifies ramp-up and drain
effects and makes multiplexing recovery more important.

Figure 1 is the model-side picture used throughout the paper. The same
figure carries both the inverted-scaling workload curves and the regime
map induced by \texttt{Re}.

Figure 1 shows the scheduling logic that the later artifact-backed
measurements are tested against.

\begin{figure}
\centering
\includegraphics[width=\textwidth,keepaspectratio,alt={Figure 1. Inverted scaling and Reynolds regime map from the Chapter 17 companion artifacts. The same analytic surface supports the speedup ceiling, the speedup floor, and the regime thresholds used to classify laminar, transitional, and turbulent execution.}]{companion-tests/artifacts/ch17-inverted-scaling-reynolds-figure.png}
\caption{Figure 1. Inverted scaling and Reynolds regime map from the
Chapter 17 companion artifacts. The same analytic surface supports the
speedup ceiling, the speedup floor, and the regime thresholds used to
classify laminar, transitional, and turbulent execution.}
\end{figure}

\subsection{3. Formal Results}\label{formal-results}

\subsubsection{3.1 Correctness Under Finite DAG
Decomposition}\label{correctness-under-finite-dag-decomposition}

The first theorem family is a sufficiency surface. Under finite DAG
decomposition and C1-C4, fork/race/fold scheduling can preserve
correctness without requiring a globally serialized path. The paper does
not claim a universal completeness theorem for all distributed systems.
It uses a narrower constructive claim: finite DAG workloads can be
decomposed into linear chains, fork points, and join points, and the
scheduler may operate on that decomposition without changing the result
when the four conditions hold.

This surface is backed by executable finite-DAG decomposition checks and
by a bounded formal stack. The TLA+ side verifies the local invariants
that define the safe scheduling envelope. The Lean side verifies the
quantitative statements used in the rest of the paper. The paper
therefore separates two questions that are often blurred together in
systems writing:

\begin{enumerate}
\def\labelenumi{\arabic{enumi}.}
\tightlist
\item
  Is the richer schedule well-defined?
\item
  Does the richer schedule help?
\end{enumerate}

The first question is answered locally by C1-C4 and the finite DAG
decomposition surface. The rest of this section answers the second
question.

\subsubsection{3.2 Pipeline Speedup Floor and
Sandwich}\label{pipeline-speedup-floor-and-sandwich}

The core scheduling theorem used here is
\texttt{THM-PIPELINE-SPEEDUP-FLOOR} from \texttt{Multiplexing.lean}.
Under A1-A2, the pipelined time never exceeds the serialized time:

\[
\lceil P/B \rceil + N - 1 \le P \cdot N.
\]

This is the lower side of the speedup sandwich. The upper side comes
from the large-\texttt{P} limit:

\[
\mathrm{speedup} \to B \cdot N \quad \text{as } P \to \infty.
\]

Together these yield the bounded interval

\[
1 \le \mathrm{speedup} \le B \cdot N.
\]

The value of the theorem is not that the ceiling is surprising. The
value is that the floor is explicit. A staged pipeline under the modeled
assumptions is not only potentially faster; it is structurally non-worse
than the serialized baseline. This is what makes a mismatch diagnostic
possible. If the workload admits a pipeline and the deployment still
measures a loss, that loss is evidence against the assumptions or
evidence that the schedule is not the schedule the model claims to
analyze.

The companion packages both the floor and the strict-speedup refinement.
The strict version applies for \texttt{P\ \textgreater{}=\ 2} and
\texttt{N\ \textgreater{}=\ 2}, ruling out the degenerate single-item
and single-stage boundaries where the pipeline collapses to the path.

\subsubsection{3.3 Reynolds
Classification}\label{reynolds-classification}

The second theorem family classifies execution regimes with the Reynolds
number. The mechanized claim used here is operational:

\begin{enumerate}
\def\labelenumi{\arabic{enumi}.}
\tightlist
\item
  \texttt{Re\ \textless{}\ 1/3} is laminar.
\item
  \texttt{1/3\ \textless{}=\ Re\ \textless{}\ 2/3} is transitional.
\item
  \texttt{Re\ \textgreater{}=\ 2/3} is turbulent.
\end{enumerate}

The companion theorem names \texttt{reynolds\_regime\_exhaustive},
\texttt{classifyRegime}, and \texttt{turbulent\_collapse\_floor} are
used because they make the schedule-level consequence explicit: the
regime is not only a label, it constrains which fold policies are safe
and how much collapse work the runtime should expect to pay.

The practical use of the classification is modest but important. It
turns the sentence ``this pipeline feels too thin to keep all stages
busy'' into a falsifiable condition on \texttt{N/C}. The paper does not
need a fluid identity claim. It only needs the regime map as a
schedule-selection tool.

\subsubsection{3.4 Delta\_beta as a Diagnostic
Coordinate}\label{delta_beta-as-a-diagnostic-coordinate}

The step-count theorem and the regime classification would already
justify a technical note on staged execution. \texttt{Delta\_beta} adds
the bridge to richer workloads. When a workload has positive intrinsic
path structure, a path-like execution collapses part of that structure
into the base path and pays for the collapse as waiting, framing, or
coordination overhead. The diagnostic claim is:

\begin{enumerate}
\def\labelenumi{\arabic{enumi}.}
\tightlist
\item
  \texttt{Delta\_beta\ =\ 0} predicts structural match.
\item
  \texttt{Delta\_beta\ \textgreater{}\ 0} predicts that some observed
  waste comes from the execution topology rather than from the problem
  alone.
\item
  Larger positive \texttt{Delta\_beta} widens the room in which a richer
  schedule can help, but does not by itself prove that a particular
  implementation will help.
\end{enumerate}

This is the point of the bridge to protocol evidence. A scheduler can
expose multiple useful paths and still lose them if the transport
collapses everything back into one coupled delivery stream. In that case
the scheduling layer and the protocol layer share the same failure: both
are enforcing a base-space path on a workload that still has unused path
structure.

The paper does not infer \texttt{Delta\_beta} from the deployed gate
scalars themselves. It computes the diagnostic on bounded graph pairs,
then asks whether the measured scheduler and transport surfaces behave
like positive-deficit cases, zero-deficit cases, or the exact one-path
boundary.

\subsubsection{3.5 Protocol Theorem
Surface}\label{protocol-theorem-surface}

The protocol story in this paper is bounded and concrete. The companion
formal surface includes the following claims.

{\def\LTcaptype{none} % do not increment counter
\begin{longtable}[]{@{}
  >{\raggedright\arraybackslash}p{(\linewidth - 4\tabcolsep) * \real{0.3333}}
  >{\raggedright\arraybackslash}p{(\linewidth - 4\tabcolsep) * \real{0.3333}}
  >{\raggedright\arraybackslash}p{(\linewidth - 4\tabcolsep) * \real{0.3333}}@{}}
\toprule\noalign{}
\begin{minipage}[b]{\linewidth}\raggedright
Claim family
\end{minipage} & \begin{minipage}[b]{\linewidth}\raggedright
Scope used here
\end{minipage} & \begin{minipage}[b]{\linewidth}\raggedright
Artifact pair
\end{minipage} \\
\midrule\noalign{}
\endhead
\bottomrule\noalign{}
\endlastfoot
quorum visibility & acknowledged values remain visible to legal quorum
reads in the bounded crash/recover model & \texttt{QuorumReadWrite.tla}
+ \texttt{QuorumVisibility.lean} \\
connected-quorum exactness & connected live quorums read committed
values exactly under explicit partition assumptions &
\texttt{QuorumAsyncNetwork.tla} + \texttt{QuorumAsyncNetwork.lean} \\
committed-session consistency & read-your-writes and monotonic reads
hold when reads are restricted to committed states &
\texttt{QuorumSessionConsistency.tla} +
\texttt{QuorumConsistency.lean} \\
multi-writer committed ordering & latest-ballot and latest-writer
properties hold under unique ordered ballots &
\texttt{QuorumMultiWriter.tla} + \texttt{QuorumOrdering.lean} \\
\end{longtable}
}

These theorems are not claims about all asynchronous distributed
protocols. They are boundary statements for the specific transport and
collapse rules used by the companion surface. Their role in the flagship
is limited: they justify saying that the protocol layer itself is not
hand-waved. The wire semantics and the fold rules are also given bounded
formal treatment.

Just as important, these theorems are not the source of the Gate2
latency numbers. In the flagship they certify bounded stream legality,
ordering, and fold semantics for the protocol surface being discussed.
Gate2 then measures a narrower question: whether Aeon Flow improves
framing and completion behavior relative to HTTP/3 on the declared
corpus and impairment matrix. Crash-recovery completeness or general
consistency claims are outside what Gate2 directly measures.

\subsection{4. Systems Instantiation}\label{systems-instantiation}

\subsubsection{4.1 Distributed Staged
Computation}\label{distributed-staged-computation}

The concrete scheduler studied in the companion is a distributed staged
computation surface for inference prefill. Prompt chunks move across
stage nodes. Each node preserves stage-local order for the stream it
owns, while the overall runtime overlaps multiple chunks across stages.
This is exactly the sort of setting where C1 matters: correctness
follows from local ordering rather than from a single total order over
all work in flight.

At the implementation level the schedule is a chunked pipeline rather
than a speculative executor. A chunk enters stage 1, advances when that
stage is free, and continues through the remaining stages while later
chunks begin behind it. The modeled makespan is the
\texttt{ceil(P/B)\ +\ N\ -\ 1} surface from Section 2.4. The measured
wall-clock evidence later in the paper tests whether the deployed system
preserves enough of that shape to separate chunked execution from the
serialized baseline.

The important structural point is that the scheduler does not treat
every token or chunk as if it needed to wait for global completion. It
exposes just enough concurrency to keep independent stage work alive
while retaining a deterministic fold at the stage boundaries. That is
the flagship example of reducing \texttt{Delta\_beta} without abandoning
determinism.

\subsubsection{4.2 Aeon Flow}\label{aeon-flow}

The scheduler runs over Aeon Flow, a binary framing layer whose unit of
transport is a \texttt{FlowFrame}. The header width is 10 bytes:

\begin{enumerate}
\def\labelenumi{\arabic{enumi}.}
\tightlist
\item
  \texttt{stream\_id}: 2 bytes
\item
  \texttt{sequence}: 4 bytes
\item
  \texttt{flags}: 1 byte
\item
  \texttt{length}: 3 bytes
\end{enumerate}

That small fact matters because the protocol question here is not only
throughput. It is whether the transport preserves enough path structure
for the scheduler to remain visible at the wire. Aeon Flow is designed
to carry multiple logical streams without forcing them through the
stronger coupling of a one-path ordered-delivery model. A stream can
therefore correspond to a computation path or a staged chunk sequence
rather than to the entire request as one inseparable line.

In the flagship argument, Aeon Flow is not interesting because it is a
custom protocol. It is interesting because it makes the scheduler's
branch structure legible at the transport layer. A richer schedule over
a transport that immediately reserializes everything would only shift
the mismatch boundary downward. The purpose of the protocol layer is to
keep the transport from becoming the hidden place where
\texttt{Delta\_beta} turns positive again.

\subsubsection{4.3 The Scheduler-Transport
Bridge}\label{the-scheduler-transport-bridge}

The scheduler and protocol bridge can be stated directly.

\begin{enumerate}
\def\labelenumi{\arabic{enumi}.}
\tightlist
\item
  The scheduler decomposes work into stage-local paths.
\item
  The transport preserves those paths as explicit streams and
  deterministic folds.
\item
  The empirical gates measure whether the resulting surface actually
  reduces observed waste.
\end{enumerate}

This is why the bridge thesis needs both gate1 and gate2, plus an
explicit transport boundary check. Gate1 checks that the staged
execution surface separates chunked and serialized wall-clock behavior
on external multi-host endpoints. The boundary checks isolate the gross
path-collapse question on bounded graph pairs. Gate2 then asks the
narrower protocol question: once the transport is at least
stream-preserving, does Aeon Flow still retain a framing and completion
advantage over HTTP/3 across a heterogeneous site corpus?

No claim in this section depends on unrelated correspondence material.
The only claim is that the runtime and the transport are shaped to
expose and preserve path structure that a path-like baseline would hide.

\subsection{5. Evaluation}\label{evaluation}

\subsubsection{5.1 Evaluation Policy}\label{evaluation-policy}

Only artifact-backed quantitative claims are used in the flagship body.
The evaluation surface is intentionally small:

\begin{enumerate}
\def\labelenumi{\arabic{enumi}.}
\tightlist
\item
  Gate1 for wall-clock staged execution.
\item
  Gate2 for protocol-corpus framing and completion behavior.
\item
  A path-boundary control showing that the same surface predicts parity
  or collapse when the path budget really is one.
\item
  The analytic figure for the Reynolds/speedup model.
\end{enumerate}

This is enough to test the bridge thesis without importing unrelated
sections of the larger chapter.

The paper only makes claims that are attached to an observable and a
failure condition.

{\def\LTcaptype{none} % do not increment counter
\begin{longtable}[]{@{}
  >{\raggedright\arraybackslash}p{(\linewidth - 6\tabcolsep) * \real{0.2500}}
  >{\raggedright\arraybackslash}p{(\linewidth - 6\tabcolsep) * \real{0.2500}}
  >{\raggedright\arraybackslash}p{(\linewidth - 6\tabcolsep) * \real{0.2500}}
  >{\raggedright\arraybackslash}p{(\linewidth - 6\tabcolsep) * \real{0.2500}}@{}}
\toprule\noalign{}
\begin{minipage}[b]{\linewidth}\raggedright
Paper-side claim
\end{minipage} & \begin{minipage}[b]{\linewidth}\raggedright
Formal or executable anchor
\end{minipage} & \begin{minipage}[b]{\linewidth}\raggedright
Observable used here
\end{minipage} & \begin{minipage}[b]{\linewidth}\raggedright
What would count against the claim
\end{minipage} \\
\midrule\noalign{}
\endhead
\bottomrule\noalign{}
\endlastfoot
chunked staged execution can remove path-imposed coordination cost under
A1-A2 & pipeline floor and sandwich theorem family & Gate1 latency
separation between serialized and chunked execution & chunked execution
fails to separate from, or reverses against, the serialized baseline on
the declared primary cells \\
\texttt{Delta\_beta} can be computed rather than only described &
bounded workload/execution graph-pair extraction plus protocol deficit
boundary checks & exact recovery of one-path, ordered-pipe-collapse, and
per-stream-preservation cases & the extracted graph pairs fail to
separate \texttt{Delta\_beta\ =\ 0} from positive-deficit boundary
cases \\
once gross path collapse has been separated, protocol structure still
changes residual coordination cost & FlowFrame protocol checks plus
Gate2 corpus harness (\texttt{Aeon\ Flow} versus \texttt{HTTP/3}) &
Gate2 framing and completion behavior across the primary environment
matrix & framing or completion gains collapse against the explicit
HTTP/3 baseline \\
\end{longtable}
}

\subsubsection{5.2 Operational Delta\_beta
Checks}\label{operational-delta_beta-checks}

Before reading the deployment gates, the paper fixes the diagnostic on
explicit graph pairs. The companion already contains the bounded cases
used for that calibration.

{\def\LTcaptype{none} % do not increment counter
\begin{longtable}[]{@{}
  >{\raggedright\arraybackslash}p{(\linewidth - 8\tabcolsep) * \real{0.1667}}
  >{\raggedleft\arraybackslash}p{(\linewidth - 8\tabcolsep) * \real{0.2222}}
  >{\raggedleft\arraybackslash}p{(\linewidth - 8\tabcolsep) * \real{0.2222}}
  >{\raggedleft\arraybackslash}p{(\linewidth - 8\tabcolsep) * \real{0.2222}}
  >{\raggedright\arraybackslash}p{(\linewidth - 8\tabcolsep) * \real{0.1667}}@{}}
\toprule\noalign{}
\begin{minipage}[b]{\linewidth}\raggedright
Case
\end{minipage} & \begin{minipage}[b]{\linewidth}\raggedleft
\texttt{beta\_1*} from workload graph
\end{minipage} & \begin{minipage}[b]{\linewidth}\raggedleft
\texttt{beta\_1\^{}exec} from effective execution graph
\end{minipage} & \begin{minipage}[b]{\linewidth}\raggedleft
\texttt{Delta\_beta}
\end{minipage} & \begin{minipage}[b]{\linewidth}\raggedright
Expected reading
\end{minipage} \\
\midrule\noalign{}
\endhead
\bottomrule\noalign{}
\endlastfoot
one-path chain & \texttt{0} & \texttt{0} & \texttt{0} & exact parity
boundary; no structural win is predicted \\
10-stream fork/join over one ordered pipe & \texttt{9} & \texttt{0} &
\texttt{9} & positive mismatch; the transport has collapsed visible
concurrency back into one path \\
10-stream fork/join over per-stream recovery (\texttt{HTTP/3} or
\texttt{Aeon\ Flow}) & \texttt{9} & \texttt{9} & \texttt{0} & stream
independence is preserved; remaining differences are protocol cost, not
gross path collapse \\
\end{longtable}
}

This is the direct operational check for the diagnostic.
\texttt{Delta\_beta} is not being reverse-engineered from Gate1 or Gate2
outcomes. It is being computed on bounded graph pairs first, then used
to sort later measurements into positive-deficit, zero-deficit, and
one-path-boundary cases.

\subsubsection{5.3 Boundary Control: The One-Path Case Must
Collapse}\label{boundary-control-the-one-path-case-must-collapse}

The same surface that supports the positive gate results also predicts a
no-win boundary. This matters because it answers the obvious reviewer
objection that the diagnostic is merely a new way to say ``parallelism
helps.'' In the companion boundary checks, setting
\texttt{beta\_1\ =\ 0} forces \texttt{C\ =\ 1}, and the pipeline
equation recovers exact sequential execution:

\[
T = (N + 1 - 1) \cdot t_{\mathrm{stage}} = N \cdot t_{\mathrm{stage}}.
\]

There is no hidden speedup term left once the path budget has collapsed
to one. The same companion boundary surface also makes the protocol
version explicit: a richer application layer can still inherit a
path-like transport bottleneck, so exposing concurrency above the wire
is not enough if the transport below it reserializes recovery and
delivery. This section is not a new headline result. It is a control. It
shows that the flagship does not use \texttt{Delta\_beta} as a universal
success score. The framework is allowed to predict parity when the
workload is one-path, and it is allowed to predict collapse when the
scheduler's extra path structure disappears at the transport layer.

\subsubsection{5.4 Gate1: External Multi-Host Wall-Clock
Evidence}\label{gate1-external-multi-host-wall-clock-evidence}

Gate1 measures chunked versus serialized execution on an external
multi-host setup. The flagship uses the external multi-host report
because it is the cleanest bridge from the formal pipeline model to
deployed staged execution. The artifact uses two workloads, multiple RTT
and loss settings, bootstrap confidence intervals, and six distinct
endpoint hosts.

The result is simple and strong. 8/8 primary Gate1 cells pass. Across
those primary cells, the median speedup range is 11.785x-21.620x. The
minimum primary-cell speedup CI low is 11.365x. Across the full ten-cell
matrix, the minimum latency-improvement CI low is 3561 ms.

Those numbers matter for two reasons. First, they are comfortably above
the speedup floor. The deployment is not merely avoiding regressions; it
is showing a large, repeated separation from the serialized baseline.
Second, the range widens as the workload gets larger in stage count and
token count, which is exactly the direction predicted by the modeled
step-count surface.

The Gate1 figure is the most direct empirical view of the scheduling
claim.

Figure 2 shows the deployed wall-clock separation that corresponds to
the model-side speedup surface.

\begin{figure}
\centering
\includegraphics[width=\textwidth,keepaspectratio,alt={Figure 2. Gate1 wall-clock evidence from the Chapter 17 companion artifacts. The chart summarizes the ten-cell external multi-host matrix and shows the separation between serialized and chunked execution together with bootstrap speedup intervals.}]{companion-tests/artifacts/ch17-gate1-wallclock-figure.png}
\caption{Figure 2. Gate1 wall-clock evidence from the Chapter 17
companion artifacts. The chart summarizes the ten-cell external
multi-host matrix and shows the separation between serialized and
chunked execution together with bootstrap speedup intervals.}
\end{figure}

The flagship does not claim that the gate1 scalar is a universal
constant. The point is narrower: on a bounded external deployment
surface, the measured behavior matches the qualitative prediction of the
formal scheduler. A path-like baseline is paying coordination cost that
the staged schedule avoids. Just as important, the one-path boundary
above tells us what would have happened if the workload or runtime had
really been \texttt{beta\_1\ =\ 0}: the measured separation would have
had no structural reason to appear.

\subsubsection{5.5 Gate2: Protocol-Corpus
Evidence}\label{gate2-protocol-corpus-evidence}

Gate2 tests the transport side of the bridge on a seeded heterogeneous
corpus. The corpus contains 144 sites and 12,371 resources, and the six
primary environments span increasing RTT, decreasing bandwidth, and
rising loss. The explicit baseline is HTTP/3. That distinction matters.
Gate2 is not where the paper establishes the positive-deficit
ordered-pipe collapse case; the companion boundary checks already do
that on bounded graph pairs. Gate2 asks the narrower question of
within-class protocol efficiency: once the transport is already
stream-preserving, do Aeon Flow's smaller framing and fold semantics
still improve framing and completion behavior relative to HTTP/3 across
the declared impairment matrix?

Again the result is compact. 6/6 primary Gate2 cells pass. The framing
median gain is 72.252\%, and the minimum primary-cell framing CI low is
72.19\%. At the benign primary edge, the \texttt{rtt12-bw80-loss0}
environment still yields 23.62 ms median completion gain and 23.79 ms
p95 completion gain. At the harsh primary edge,
\texttt{rtt110-bw7-loss2pct} yields 92.39 ms median completion gain and
110.5 ms p95 completion gain.

This is the key within-class transport result. The framing advantage is
nearly invariant across environments because it comes from protocol
structure. The completion gain grows with impairment because the lower
framing and recovery burden compounds as the network gets worse. The
flagship does not read Gate2 as proof that HTTP/3 is a one-path
transport. It reads Gate2 more narrowly: after the path-collapse
boundary has been separated elsewhere, protocol structure still changes
the residual coordination cost of carrying the staged workload over the
wire.

Figure 3 makes the split visible: nearly constant framing advantage,
widening completion advantage.

\begin{figure}
\centering
\includegraphics[width=\textwidth,keepaspectratio,alt={Figure 3. Gate2 protocol-corpus evidence from the Chapter 17 companion artifacts. The chart separates framing gains from completion gains across the heterogeneous site corpus and shows that the protocol advantage persists across all six primary environments.}]{companion-tests/artifacts/ch17-gate2-protocol-corpus-figure.png}
\caption{Figure 3. Gate2 protocol-corpus evidence from the Chapter 17
companion artifacts. The chart separates framing gains from completion
gains across the heterogeneous site corpus and shows that the protocol
advantage persists across all six primary environments.}
\end{figure}

\subsubsection{5.6 Interpreting the Evaluation Surface
Together}\label{interpreting-the-evaluation-surface-together}

The graph-pair checks, Gate1, and Gate2 do different jobs.

\begin{enumerate}
\def\labelenumi{\arabic{enumi}.}
\tightlist
\item
  The graph-pair checks operationalize \texttt{Delta\_beta} and separate
  one-path, ordered-pipe-collapse, and per-stream-preserving cases
  before any deployment scalar is discussed.
\item
  Gate1 checks the staged scheduler against the serialized baseline.
\item
  Gate2 checks Aeon Flow against the explicit HTTP/3 baseline once the
  gross collapse-vs-preservation question has already been isolated.
\item
  The one-path control checks that the same formal surface also predicts
  the no-gain case when the path budget really is one.
\end{enumerate}

Taken together they support one bounded conclusion: bounded graph pairs
make \texttt{Delta\_beta} operational enough to separate
positive-deficit path-collapse cases from matched cases, and on that
foundation the staged runtime plus transport surface shows
artifact-backed reductions in observed coordination cost. That is the
precise role of \texttt{Delta\_beta} in the flagship: not a one-number
oracle, but a way to state where missing path budget is being inserted.

The paper does not need more than that. It does not need a claim that
every transport problem is topological, or that every pipeline problem
should be reframed through Betti numbers. It only needs the evidence to
show that the diagnostic is usable on the targeted
scheduler-plus-transport bridge, that it cleanly recovers the one-path
and ordered-pipe boundaries, and that the deployment gates do not
contradict those local predictions.

\subsection{6. Limitations and
Non-Claims}\label{limitations-and-non-claims}

The flagship deliberately narrows scope. The following boundaries are
part of the claim, not footnotes to it.

\begin{enumerate}
\def\labelenumi{\arabic{enumi}.}
\tightlist
\item
  The model is a finite DAG model. Infinite-state, adversarially
  adaptive, or unrestricted asynchronous systems are outside the direct
  theorem surface used here.
\item
  The pipeline speedup formula is a modeled step-count claim under
  A1-A2. It is not a universal wall-clock predictor.
\item
  \texttt{Delta\_beta} is a diagnostic coordinate, not a one-number
  oracle for system quality. A positive value says the execution
  topology is narrower than the modeled workload, not that a particular
  implementation must improve by a fixed amount.
\item
  Gate1 is a bounded external multi-host matrix, not a claim about every
  deployment.
\item
  Gate2 is a seeded heterogeneous corpus with explicit impairment
  environments, and its protocol comparison is Aeon Flow versus HTTP/3
  rather than a universal transport shootout.
\item
  The paper does not claim universal superiority of richer schedules.
  For small workloads, extremely fast serial paths, or high coordination
  overhead, the correct topology may be \texttt{beta\_1\ =\ 0}.
\item
  No body claim in this paper depends on extra-domain correspondence
  surfaces outside distributed systems and bounded protocol models.
\end{enumerate}

These limitations sharpen the result rather than weaken it. The claim is
strong precisely because it stays local.

\subsection{7. Related Work}\label{related-work}

This section is intentionally not a survey. It covers only the
literatures needed to identify the incremental claim of the flagship.

The speedup formulas in this paper sit beside, not above, the classical
scaling laws. Amdahl studies fixed-workload speed limits under an
explicit serial fraction, while Gustafson reframes the bound under
scaled workloads {[}1, 2{]}. Those laws remain the right tools for
ceiling questions. The current paper adds a different question: when is
part of the apparent serial fraction coming from the execution surface
rather than from the workload's natural structure? \texttt{Delta\_beta}
is meant as a local mismatch diagnostic for that narrower question, not
as a replacement for classical scaling analysis.

The scheduling model is likewise adjacent to older work on explicit
concurrency surfaces. Kahn process networks and early dataflow
architectures made it normal to reason in terms of bounded dependency
graphs rather than one total order {[}3, 4{]}. MapReduce made
fork-and-fold structure operational on large bounded computations
{[}6{]}. The flagship does not claim novelty for DAG execution itself.
Its narrower claim is that separating workload-side path demand from
runtime-side path budget makes certain scheduling failures easier to
localize, especially when the system appears parallel at one layer and
path-like at another.

Classical pipeline overlap and asynchronous stage composition are also
established. Micropipelines already show that local stage ordering can
be enough to preserve correctness while allowing overlap {[}5{]}. The
scheduler claim here is therefore not ``pipelining exists.'' It is that
a staged system can still be under-provisioned in path structure
relative to the workload it carries, and that the under-provisioning can
be stated as a path-budget mismatch rather than only as low occupancy or
poor hardware utilization.

The protocol side belongs to a familiar lineage as well. HTTP/2 reduces
framing waste and introduces application-layer multiplexing, while QUIC
moves recovery and ordering decisions closer to individual streams {[}8,
9{]}. The paper does not claim to supersede that literature. Its
narrower bridge claim is that transport should be evaluated partly by
whether it preserves the scheduler's exposed path structure or silently
collapses it back toward one ordered recovery path. That is why the
flagship studies scheduler and protocol together instead of treating
them as unrelated optimizations.

Finally, the formal-methods layer is deliberately conventional in the
best sense. TLA+ and related specification practice already provide a
language for bounded protocol safety and progress arguments {[}7{]}. The
novelty claim here is not ``formal methods for systems'' in general. It
is the tighter connection between a bounded formal surface, a small set
of measured observables, and explicit falsifiers in the main paper, all
backed by a public rerun surface rather than by prose alone {[}10{]}.
That is what is meant to keep the flagship from collapsing into either
pure systems benchmarking or pure notation.

\subsection{8. Conclusion}\label{conclusion}

Path-like execution topologies impose avoidable coordination cost when
the workload's natural structure has positive parallel-path complexity.

That sentence is the whole paper. The finite DAG model provides the
vocabulary. C1-C4 provide the correctness envelope. The speedup floor,
speedup sandwich, and Reynolds classification provide the model-side
quantitative surface. Aeon Flow provides a transport that preserves the
scheduler's path structure instead of silently collapsing it. Gate1 and
Gate2 show that the bridge is not only formal: the staged runtime and
the protocol layer both exhibit artifact-backed gains on bounded
evaluation surfaces.

The result is narrow, but it is enough. One does not need a universal
theory of computation to say something useful about topological mismatch
in distributed inference. One needs a finite DAG model, explicit
assumptions, mechanized local claims, and measured evidence. That is
what this flagship paper provides.

The topological vocabulary only matters to the extent that it
distinguishes two things standard prose often blends together: path
structure inherent to the workload and path collapse inserted by the
execution surface. That is the only place \texttt{Delta\_beta} is meant
to earn its keep.

\subsection{Appendix A: Theorem-to-Artifact
Map}\label{appendix-a-theorem-to-artifact-map}

{\def\LTcaptype{none} % do not increment counter
\begin{longtable}[]{@{}
  >{\raggedright\arraybackslash}p{(\linewidth - 4\tabcolsep) * \real{0.3333}}
  >{\raggedright\arraybackslash}p{(\linewidth - 4\tabcolsep) * \real{0.3333}}
  >{\raggedright\arraybackslash}p{(\linewidth - 4\tabcolsep) * \real{0.3333}}@{}}
\toprule\noalign{}
\begin{minipage}[b]{\linewidth}\raggedright
Local claim in flagship
\end{minipage} & \begin{minipage}[b]{\linewidth}\raggedright
Theorem or artifact family
\end{minipage} & \begin{minipage}[b]{\linewidth}\raggedright
Concrete artifact
\end{minipage} \\
\midrule\noalign{}
\endhead
\bottomrule\noalign{}
\endlastfoot
finite DAG scheduling under explicit decomposition assumptions &
executable finite-DAG decomposition surface &
\texttt{companion-\allowbreak{}tests/\allowbreak{}src/\allowbreak{}dag-\allowbreak{}completeness.\allowbreak{}test.\allowbreak{}ts} \\
C1-C4 correctness envelope & bounded TLA+ and Lean scheduler surface &
\texttt{TemporalModel} companion modules plus \texttt{@a0n/aeon-logic}
rerun path \\
pipeline floor \texttt{T\_pipe\ \textless{}=\ T\_seq} &
\texttt{THM-PIPELINE-SPEEDUP-FLOOR} &
\texttt{formal/lean/Multiplexing.lean} \\
speedup interval \texttt{1\ \textless{}=\ speedup\ \textless{}=\ B*N} &
\texttt{pipeline\_speedup\_floor}, \texttt{pipeline\_strict\_speedup},
\texttt{pipeline\_speedup\_sandwich} &
\texttt{formal/lean/Multiplexing.lean},
\texttt{formal/lean/PredictionsRound15.lean} \\
Reynolds regime cut points & \texttt{reynolds\_regime\_exhaustive},
\texttt{classifyRegime}, \texttt{turbulent\_collapse\_floor} &
\texttt{formal/lean/ReynoldsBFT.lean},
\texttt{companion-\allowbreak{}tests/\allowbreak{}artifacts/\allowbreak{}ch17-\allowbreak{}inverted-\allowbreak{}scaling-\allowbreak{}reynolds-\allowbreak{}figure.\allowbreak{}\{json,md,png,svg\}} \\
operational \texttt{Delta\_beta} extraction from bounded graph pairs &
workload/execution graph comparison over explicit traces &
\texttt{companion-\allowbreak{}tests/\allowbreak{}src/\allowbreak{}deficit-\allowbreak{}evidence.\allowbreak{}test.\allowbreak{}ts} \\
10-byte FlowFrame header and stream semantics & framing law plus runtime
protocol tests & \texttt{companion-\allowbreak{}tests/\allowbreak{}src/\allowbreak{}flow-\allowbreak{}protocol.\allowbreak{}test.\allowbreak{}ts} \\
bounded quorum visibility and ordering surface & quorum theorem family &
\texttt{QuorumReadWrite.tla}, \texttt{QuorumVisibility.lean},
\texttt{QuorumAsyncNetwork.tla}, \texttt{QuorumAsyncNetwork.lean},
\texttt{QuorumSessionConsistency.tla}, \texttt{QuorumConsistency.lean},
\texttt{QuorumMultiWriter.tla}, \texttt{QuorumOrdering.lean} \\
one-path boundary recovers sequential execution & explicit
\texttt{beta\_1\ =\ 0} boundary check &
\texttt{companion-\allowbreak{}tests/\allowbreak{}src/\allowbreak{}pipeline-\allowbreak{}topology.\allowbreak{}test.\allowbreak{}ts} \\
transport can reintroduce path collapse below the application layer &
protocol deficit boundary &
\texttt{companion-\allowbreak{}tests/\allowbreak{}src/\allowbreak{}deficit-\allowbreak{}evidence.\allowbreak{}test.\allowbreak{}ts} \\
deployed scheduler evidence & Gate1 wall-clock matrix &
\texttt{companion-\allowbreak{}tests/\allowbreak{}artifacts/\allowbreak{}gate1-\allowbreak{}wallclock-\allowbreak{}external-\allowbreak{}multihost.\allowbreak{}json},
\texttt{companion-\allowbreak{}tests/\allowbreak{}artifacts/\allowbreak{}ch17-\allowbreak{}gate1-\allowbreak{}wallclock-\allowbreak{}figure.\allowbreak{}\{json,md,png,svg\}} \\
deployed protocol evidence & Gate2 protocol corpus &
\texttt{companion-\allowbreak{}tests/\allowbreak{}artifacts/\allowbreak{}gate2-\allowbreak{}protocol-\allowbreak{}corpus.\allowbreak{}json},
\texttt{companion-\allowbreak{}tests/\allowbreak{}artifacts/\allowbreak{}ch17-\allowbreak{}gate2-\allowbreak{}protocol-\allowbreak{}corpus-\allowbreak{}figure.\allowbreak{}\{json,md,png,svg\}} \\
\end{longtable}
}

\subsection{Appendix B: Reproduction
Surface}\label{appendix-b-reproduction-surface}

The flagship paper has its own build and validation path so that it can
be checked independently from the full theory chapter.

\subsubsection{Build}\label{build}

\begin{Shaded}
\begin{Highlighting}[]
\ExtensionTok{pnpm}\NormalTok{ run manuscript:arxiv:flagship:tex}
\ExtensionTok{pnpm}\NormalTok{ run manuscript:arxiv:flagship:build}
\ExtensionTok{pnpm}\NormalTok{ run manuscript:arxiv:flagship:package}
\end{Highlighting}
\end{Shaded}

\subsubsection{Targeted manuscript
validation}\label{targeted-manuscript-validation}

\begin{Shaded}
\begin{Highlighting}[]
\ExtensionTok{pnpm} \AttributeTok{{-}{-}dir}\NormalTok{ open{-}source/aeon/docs/ebooks/145{-}log{-}rolling{-}pipelined{-}prefill/companion{-}tests exec vitest run }\DataTypeTok{\textbackslash{}}
\NormalTok{  src/manuscript{-}flagship{-}structure.test.ts }\DataTypeTok{\textbackslash{}}
\NormalTok{  src/manuscript{-}flagship{-}artifact{-}consistency.test.ts }\DataTypeTok{\textbackslash{}}
\NormalTok{  src/manuscript{-}flagship{-}hardening.test.ts}
\end{Highlighting}
\end{Shaded}

\subsubsection{Supporting scheduler and protocol
checks}\label{supporting-scheduler-and-protocol-checks}

\begin{Shaded}
\begin{Highlighting}[]
\ExtensionTok{pnpm} \AttributeTok{{-}{-}dir}\NormalTok{ open{-}source/aeon/docs/ebooks/145{-}log{-}rolling{-}pipelined{-}prefill/companion{-}tests exec vitest run }\DataTypeTok{\textbackslash{}}
\NormalTok{  src/dag{-}completeness.test.ts }\DataTypeTok{\textbackslash{}}
\NormalTok{  src/pipeline{-}formulas.test.ts }\DataTypeTok{\textbackslash{}}
\NormalTok{  src/pipeline{-}topology.test.ts }\DataTypeTok{\textbackslash{}}
\NormalTok{  src/deficit{-}evidence.test.ts }\DataTypeTok{\textbackslash{}}
\NormalTok{  src/flow{-}protocol.test.ts }\DataTypeTok{\textbackslash{}}
\NormalTok{  src/ch17{-}gate1{-}wallclock{-}figure.test.ts }\DataTypeTok{\textbackslash{}}
\NormalTok{  src/ch17{-}gate2{-}protocol{-}corpus{-}figure.test.ts}
\end{Highlighting}
\end{Shaded}

\subsubsection{Packaged submission
helper}\label{packaged-submission-helper}

\begin{Shaded}
\begin{Highlighting}[]
\FunctionTok{bash}\NormalTok{ open{-}source/aeon/docs/ebooks/145{-}log{-}rolling{-}pipelined{-}prefill/arxiv{-}submission{-}flagship/prepare{-}arxiv.sh}
\end{Highlighting}
\end{Shaded}

The flagship uses the same checked-in companion artifact directory as
the full chapter. It does not define a second artifact format. The
difference is only the manuscript selection, the narrowed claim surface,
and the dedicated build and packaging entrypoints.

\subsection{References}\label{references}

{[}1{]} G. M. Amdahl, ``Validity of the Single Processor Approach to
Achieving Large-Scale Computing Capabilities,'' \emph{AFIPS Spring Joint
Computer Conference}, 30:483-485, 1967.

{[}2{]} J. L. Gustafson, ``Reevaluating Amdahl's Law,''
\emph{Communications of the ACM}, 31(5):532-533, 1988.

{[}3{]} G. Kahn, ``The Semantics of a Simple Language for Parallel
Programming,'' \emph{IFIP Congress}, 1974.

{[}4{]} J. B. Dennis, D. P. Misunas, ``A Preliminary Architecture for a
Basic Data-Flow Processor,'' \emph{Proceedings of the 2nd Annual
Symposium on Computer Architecture}, 1975.

{[}5{]} I. E. Sutherland, ``Micropipelines,'' \emph{Communications of
the ACM}, 32(6):720-738, 1989.

{[}6{]} J. Dean, S. Ghemawat, ``MapReduce: Simplified Data Processing on
Large Clusters,'' \emph{OSDI}, 2004.

{[}7{]} L. Lamport, \emph{Specifying Systems: The TLA+ Language and
Tools for Hardware and Software Engineers}, Addison-Wesley, 2002.

{[}8{]} M. Belshe, R. Peon, M. Thomson, ``Hypertext Transfer Protocol
Version 2 (HTTP/2),'' RFC 7540, 2015.

{[}9{]} J. Iyengar, M. Thomson, ``QUIC: A UDP-Based Multiplexed and
Secure Transport,'' RFC 9000, 2021.

{[}10{]} T. W. Buley, ``Aeon Clockwork: A Unified Probability Engine
with Immanent Self-Verification,'' open-source implementation, 2026.
\url{https://github.com/forkjoin-ai/aeon-clockwork}

\end{document}
